\section{Round 2 : voisins et indicateurs sans labels}\label{sec:voisins}

Ces outils ont été introduits en v0.2.0 (round 2), puis exécutés sur les vraies données en V3 : le vote entre voisins y fait passer le LB de $0{,}6028$ à $\mathbf{0{,}6351}$ (section~\ref{sec:v3}). Il reste dans toutes les soumissions suivantes, et sert aussi à construire le professeur de l'auto-apprentissage (section~\ref{sec:selftrain}). Cette section décrit la méthode ; les résultats sont à la section~\ref{sec:resultats}.

\subsection{Faire voter les fenêtres d'un même titre-jour}
Le test contient environ 20 fenêtres par titre et par jour, et elles partagent l'effet $\eta_{c,d}$ de~\eqref{eq:hier}. Le mécanisme qui gonflait la validation aléatoire (reconnaître le jour) peut être retourné en notre faveur : si une représentation $z$ rapproche les fenêtres d'un même titre-jour, une fenêtre ambiguë peut être corrigée par ses voisines plus nettes.

\paragraph{Graphe des $k$ plus proches voisins.} Pour un ensemble de fenêtres (le test entier, ou valid $\cup$ stress pour l'évaluation), on normalise $\hat z_i=z_i/\|z_i\|$, on note $\mathcal{N}_k(i)$ les $k$ indices maximisant $\hat z_i^\top\hat z_j$ avec $j\neq i$, et l'on forme $A_{ij}=\frac1k\ind{j\in\mathcal{N}_k(i)}$. Le calcul est exact, par blocs de 4\,096 lignes, sur GPU si disponible.

\paragraph{Propagation.} Pour $\alpha\in[0,1)$ :
\begin{equation}
Q^{(0)}=P,\qquad Q^{(n+1)}=(1-\alpha)P+\alpha AQ^{(n)} .
\end{equation}
\begin{proposition}
Pour $\alpha<1$, $Q^{(n)}\to Q^{(\infty)}=(1-\alpha)(I-\alpha A)^{-1}P$.
\end{proposition}
\begin{proof}
$A$ est stochastique par lignes, donc $\|A\|_\infty=1$ et $\|\alpha A\|_\infty<1$. L'itération est contractante, de point fixe unique $Q=(1-\alpha)P+\alpha AQ$. La série de Neumann $\sum_n(\alpha A)^n$ converge vers $(I-\alpha A)^{-1}$.
\end{proof}
\noindent C'est la propagation d'étiquettes de \cite{zhou}, appliquée à des probabilités \emph{prédites} plutôt qu'à des labels connus. Le code utilise 1 ou 5 itérations, puis renormalise les lignes et applique Sinkhorn.

\paragraph{Trois représentations.}
\begin{itemize}
\item \code{window} : blocs \code{levels}, \code{ticks}, \code{lots}, \code{cat\_freq} et \code{position}, standardisés à l'intérieur de l'ensemble. Ce sont justement des statistiques de régime : leur dérive nuit à la classification, mais elles identifient bien le jour.
\item \code{embedding} : $[z_{\text{seq}}\|z_{\text{ctx}}]$ des modèles, normalisé puis moyenné entre les membres. Les modèles de développement servent pour valid/stress, ceux du refit pour le test. Le code refuse d'utiliser les poids du refit sur des partitions étiquetées.
\item \code{probs} : $\sqrt{P}$, soit une distance de Hellinger.
\end{itemize}

\paragraph{Pureté.} Avant tout lissage, on mesure la part des voisins qui appartiennent au même titre :
\begin{equation}
\text{pureté}@k=\frac1{nk}\sum_i\sum_{j\in\mathcal{N}_k(i)}\ind{y_j=y_i}.
\end{equation}

\paragraph{Densité et facteur $k$.} Le pool d'évaluation valid $\cup$ stress contient $25\,\%$ des fenêtres de chaque jour, soit environ 5 par titre-jour, alors que le test en contient 20. Le $k$ choisi sur le pool est donc multiplié par $4$ pour le test. C'est une \emph{hypothèse}, écrite dans le code (\code{k\_scale}).

\paragraph{Règle d'adoption.} Le lissage n'est adopté que s'il améliore la précision \emph{équilibrée} sur le pool d'au moins $0{,}5$ point. Ce gain est un score de sélection, pas une estimation du gain sur le test.

\subsection{Test préliminaire : l'espace des probabilités ne suffit pas}\label{sec:probaspace}
Sur les probabilités réelles du run 1, avec des voisins dans l'espace \code{probs} :
\begin{center}\small
\begin{tabular}{lccc}
\toprule
pool valid $\cup$ stress & brut & équilibré & meilleur lissage + équilibrage\\
\midrule
précision & 0,8079 & 0,8281 & 0,8281\\
\bottomrule
\end{tabular}
\end{center}
C'était attendu : deux fenêtres voisines \emph{en probabilités} portent la même information, donc la propagation n'ajoute rien. L'intérêt du lissage dépend d'une représentation qui encode le \emph{jour}, ce que le round 2 mesure avec la pureté.

\subsection{Indicateurs sans labels sur le test}
Pour des probabilités test $P$ (modèle de développement) :
\begin{align}
\text{confiance}&=\tfrac1N\textstyle\sum_i\max_kP_{ik},\qquad
\text{entropie}=-\tfrac1N\sum_{i,k}P_{ik}\log P_{ik},\\
\text{KL d'équilibre}&=\textstyle\sum_k\hat\pi_k\log(K\hat\pi_k),\quad \hat\pi_k=\tfrac1N\sum_i\ind{\arg\max P_i=k},\\
\text{accord entre graines}&=\text{accord}(m,m')\ \text{pour deux graines d'une même configuration.}
\end{align}
Au run 1, pour \code{signature\_s0} :
\begin{center}\small
\begin{tabular}{lcc}
\toprule
& valid & test\\
\midrule
confiance & 0,858 & 0,738\\
accord entre graines & 0,845 & 0,648\\
plus grand / plus petit nombre de prédictions par titre & $\approx1$ & 5,4\\
\bottomrule
\end{tabular}
\end{center}
Ces chiffres annonçaient, avant toute soumission, que le test serait bien plus difficile que la validation. Ils ne sont pas des accuracies. Ils ne servent à présélectionner des modèles qu'une fois \emph{ancrés} à quelques scores LB connus.
